\documentclass[12pt]{article}

\setlength{\parindent}{0pt}
\setcounter{secnumdepth}{3}
\usepackage[margin=1in]{geometry}

% packages
\usepackage{amsmath, amsfonts, amsthm, amssymb}
\usepackage{graphicx, float}
\usepackage{mathtools}
\usepackage{titlesec}
\usepackage{interval}
\usepackage{hyperref}
\usepackage{siunitx}
\usepackage{titling}
\usepackage{vwcol}
\usepackage{setspace}
\usepackage{empheq}
\usepackage{cancel}
\usepackage{esdiff}
\usepackage{multicol}
\usepackage{mdframed}
\usepackage{esdiff}
\usepackage{tikzsymbols}
\usepackage{multicol}
\usepackage{tikz}
\usepackage{varwidth}
\usepackage{pgfplots}
\usepackage{fancyhdr}
\usepackage{enumitem}
\usepackage{tcolorbox}

\intervalconfig {
	soft open fences
}

% header
\pagestyle{fancy}
\fancyhf{}
\lhead{Ethan Fan}
\rhead{AP Calculus BC}
\cfoot{\thepage}

% section format
\titleformat{\section}
  {\large \bfseries}
  {}
  {0em}
  {}
\titlespacing*{\section}{0pt}{0.5em}{0.3em}

\titleformat{\subsection}[runin]
  {\bfseries}
  {}
  {0em}
  {}
\titlespacing*{\subsection}{0pt}{0pt}{0em}

\titleformat{\subsubsection}[runin]
  {\itshape}
  {(\roman{subsubsection})}
  {0em}
  {}

% commands
\newcommand{\alignedintertext}[1]{%
  \noalign{%
    \vskip\belowdisplayshortskip
    \vtop{\hsize=\linewidth#1\par
    \expandafter}%
    \expandafter\prevdepth\the\prevdepth
  }%
}

\newcommand*{\paren}[1]{\ensuremath\left(#1\right)}
\newcommand*{\limit}[2][x]{\ensuremath{\displaystyle\lim_{#1\to#2}}}
\newcommand*{\Limit}[3][x]{\ensuremath{\displaystyle\lim_{#1\to#2}\left[#3\right]}}
\newcommand*{\deriv}[1][x]{\ensuremath{\dfrac{\mathrm{d}}{\mathrm{d}#1}}}
\newcommand*{\Deriv}[2][x]{\ensuremath{\dfrac{\mathrm{d}}{\mathrm{d}#1}\left[#2\right]}}
\newcommand{\ddx}{\dfrac{d}{dx}}
\newcommand{\dydx}{\dfrac{dy}{dx}}
\newcommand{\dydt}{\dfrac{dy}{dt}}
\newcommand{\dxdt}{\dfrac{dx}{dt}}
\newcommand*{\iinteg}[2][x]{\ensuremath{\displaystyle\int #2\;\mathrm{d}#1}}
\newcommand*{\dinteg}[4][x]{\ensuremath{\displaystyle\int_{#2}^{#3}#4\;\mathrm{d}#1}}
\newcommand*{\abs}[1]{\ensuremath{\left|#1\right|}}
\newcommand*{\eps}{\varepsilon}
\newcommand*{\real}{\ensuremath\mathbb{R}}
\newcommand*{\integer}{\ensuremath\mathbb{Z}}
\newcommand*{\posint}{\ensuremath\mathbb{Z^+}}
\newcommand*{\cbrt}[1]{\ensuremath{\sqrt[3]{#1}}}
\newcommand*{\lheqzero}{\ensuremath{\underset{\text{L'H}}{\overset{\left[\frac00\right]}{=}}}}
\newcommand*{\lheqinfty}{\ensuremath{\underset{\text{L'H}}{\overset{\left[\frac{\infty}{\infty}\right]}{=}}}}
\newcommand{\tor}{\text{ or }}
\newcommand{\floor}[1]{\lfloor #1 \rfloor}

\newcommand{\vem}{\vspace{0.5em}}
\newcommand{\example}[1]{\section{Example #1}}
\newcommand{\problem}[1]{\section{Problem #1}}
\newcommand{\subproblem}[1]{\subsection{(#1)} \,}

\DeclareMathOperator{\dne}{DNE}
\DeclareMathOperator{\sgn}{sgn}

\newcommand{\definition}[1]{\begin{tcolorbox}[colback=red!5!white,colframe=red!75!black,parbox=false] #1 \end{tcolorbox}}

\newcommand{\theorem}[2]{\begin{tcolorbox}[title={#1},colback=blue!5!white,colframe=blue!75!black,parbox=false] #2 \end{tcolorbox}}

\allowdisplaybreaks
% \postdisplaypenalty=100000

% document
\begin{document}

\vspace*{0cm}

% opening
\begin{center}
    {\Large \bfseries Problem Set 12} \\[0.5cm]
    {\large Odd, Even, and Periodic Functions} \\[0.35cm]
    {\normalsize October 1, 2026}
\end{center}

% problems

\problem{2}

\subproblem{a}
\begin{gather*}
    f(x)=-f(-x),\, g(x)=-g(-x)\\
    h(x)=f(x)g(x)\\
    h(-x)=f(-x)g(-x)=(-f(x))(-g(x))=h(x)\\
    \boxed{h(x) \text{ is even}}\\
    f(x)=-f(-x),\, h(x)=h(-x)\\
    k(x)=f(x)h(x)\\
    k(-x)= f(-x)h(-x)=-f(x)h(x)=-k(x)\\
    \boxed{k(x) \text{ is odd}}
\end{gather*}

\subproblem{b}
\begin{gather*}
    f(x)=f(-x)\\
    f'(x)=-f'(-x)\\
    \boxed{f'(x)\text{ is odd}}\\
    g(x)=-g(-x)\\
    g'(x)=g'(-x)\\
    \boxed{g'(x)\text{ is even}}
\end{gather*}

\subproblem{c}
\begin{gather*}
    f(x)=-f(-x)\\
    g(x)=\int_a^x f(t)\, dt\\
    g'(x)=f(x)\\
    h(x)=g(x)-g(-x)\\
    h'(x)=g'(x)+g'(-x)=f(x)+f(-x)=0\\
    h(0)=g(0)-g(0)=0\\
    g(x)-g(-x)=0 \implies \boxed{g(x)=g(-x)}
\end{gather*}

\subproblem{d}
\begin{gather*}
    a=0 \implies F(x)=\int_0^x f(t)\, dt\\
    f(t)=f(-t)\\
    F(-x)=\int_0^{-x} f(t)\, dt= -\int_0^x f(-x)\, dx = -\int_0^x f(x)\, dx \\
    \boxed{F(x)=-F(-x)}\\
    a\ne 0,\, f(x)=1\\
    F(x)=\int_1^x f(t)\, dt = x-1\\
    F(-x) = -x-1\\
    \boxed{F(x)\ne -F(-x)}
\end{gather*}

\subproblem{e}
No, only one where the constant $C=0$ is odd.

\problem{3}

\subproblem{d}
\begin{align*}
    &\int_{-2}^2 \frac{\arcsin (\sin x)+\arccos (\cos x)}{(1+x^2)\paren{1+\floor{\frac{x^2}{5}}}}\, dx\\
    =&\int_{-2}^2 \frac{\arcsin (\sin x)}{(1+x^2)\paren{1+\floor{\frac{x^2}{5}}}}\, dx + \int_{-2}^2 \frac{\arccos (\cos x)}{(1+x^2)\paren{1+\floor{\frac{x^2}{5}}}}\, dx\\
    =& \, 0+ 2 \int_0^2 \frac{\arccos (\cos x)}{(1+x^2)\paren{1+\floor{\frac{x^2}{5}}}}\, dx\\
    =&\, 2 \int_0^2 \frac{x}{1+x^2} \, dx\\
    = &\, \int_1^5 \frac{1}{u} \, du\\
    =&\, \ln 5 - \ln 1\\
    =&\,\boxed{\ln 5}
\end{align*}

\problem{4}

\subproblem{a}
\begin{gather*}
    E(x)=\frac{f(x)+f(-x)}{2}\\
    E(-x)=\frac{f(-x)+f(-(-x))}{2}=\frac{f(x)+f(-x)}{2}\\
    \boxed{E(x)=E(-x)}\\
    O(x)=\frac{f(x)-f(-x)}{2}\\
    O(-x)=\frac{f(-x)-f(-(-x))}{2}=\frac{f(-x)-f(x)}{2}\\
    \boxed{O(x)=-O(-x)}\\
    f(x)=\frac{f(x)+f(-x)}{2}+\frac{f(x)-f(-x)}{2}=E(x)+O(x)
\end{gather*}

\subproblem{b}
\begin{gather*}
    \cosh x=\frac{e^x+e^{-x}}{2}\\
    \sinh x=\frac{e^x-e^{-x}}{2}
\end{gather*}

\subproblem{c}
\begin{gather*}
    \int_{-a}^af(x)\, dx=\int_0^af(x)\, dx+\int_{-a}^0 f(x)\, dx=\int_0^a[f(x)+f(-x)]\, dx
\end{gather*}

\subproblem{d}
\begin{gather*}
    \frac{1}{1+e^x}+\frac{1}{1+e^{-x}}=\frac{2+e^x+e^{-x}}{2+e^x+e^{-x}}=1\\
    \int_{-1}^1\frac{dx}{1+e^x} = \int_0^1 \frac{1}{1+e^x}+\frac{1}{1+e^{-x}}\, dx= \int_0^1 1\, dx=\boxed{1}
\end{gather*}

\subproblem{e}
\begin{gather*}
    \frac{x^2}{1+e^{x^3}}+\frac{x^2}{1+e^{-x^3}} = x^2\\
    \int_{-2}^2 \frac{x^2}{1+e^{x^3}}\, dx = \int_0^2 \frac{x^2}{1+e^{x^3}}+\frac{x^2}{1+e^{-x^3}} \, dx = \int_0^2 x^2\,dx = \boxed{\frac{8}{3}}
\end{gather*}

\problem{5}

\subproblem{a}
\begin{gather*}
    \int_a^{a+T} f(x)\, dx= \int_a^T f(x)\, dx+\int_T^{a+T} f(x)\, dx\\
    = \int_a^T f(x)\, dx+\int_0^{a} f(x)\, dx = \int_0^T f(x)\, dx
\end{gather*}

\subproblem{b}
\begin{gather*}
    \int_0^{nT}f(x)\, dx = \sum_{k=0}^{n-1} \int_{kT}^{(k+1)T} f(x)\, dx = n\int_0^T f(x)\,dx
\end{gather*}

\subproblem{c}
\begin{gather*}
    T_{|\sin x|} = \pi \\
    \int_0^\frac{16\pi}{3} |\sin x|\, dx =  5\int_0^\pi |\sin x| \, dx + \int_0^\frac{\pi}{3}|\sin x |\, dx\\
    =5\cdot 2+1-\frac{1}{2}=\boxed{\frac{21}{2}}
\end{gather*}

\subproblem{d}
\begin{gather*}
    F(x+T)-F(x)=\int_0^{x+T} f(t)\, dt - \int_0^x f(t)\, dt=\int_x^{x+T} f(t)\, dt = \int_0^T f(t)\, dt\\
    F(x+T) = F(x) \implies F(x+T)-F(x)=0\\
    \iff  \int_0^T f(t)\, dt =0
\end{gather*}

\subproblem{e}
As $\int_0^T f(t)\, dt \ne 0$, the shape of $F(x)$ is not periodic.

\problem{7}

\subproblem{a}
\begin{gather*}
    t\in \{0,2\pi,4\pi\} \implies y=1\\
    t\in (0, \frac{\pi}{2}] \cup  [\frac{3\pi}{2}, 2\pi )\cup (2\pi ,  \frac{5\pi}{2}] \cup [\frac{7\pi}{2},4\pi) \implies y=0\\
    t\in (\frac{\pi}{2}, \frac{3\pi}{2}) \cup (\frac{5\pi}{2},\frac{7\pi}{2})\implies y=-1\\
    T=2\pi 
\end{gather*}

\subproblem{b}
\begin{gather*}
    \int_0^{2\pi} \floor{\cos x} = -\pi 
\end{gather*}
The isolated point does not matter, as its width is too insignificant compared to the rest of the function. \vem

\subproblem{c}
\begin{gather*}
    \int_0^x \floor{\cos t}\, dt = \int_0^{2\pi n} \floor{\cos t}\, dt = n\cdot -\pi =\boxed{-n\pi}
\end{gather*}

\subproblem{d}
G is constant on the intervals where $y=0$ and decreasing on the intervals where $y=-1$.\vem

\subproblem{e}
\begin{gather*}
    T=1\\
    \int_0^{7.5} s(t)\,d t=7\cdot \frac{1}{2}+ \frac{1}{2}\cdot \frac{1}{4}=\boxed{\frac{29}{8}}
\end{gather*}

\problem{8}

\subproblem{a}
\begin{gather*}
    \frac{g(x)}{1+b^x}+\frac{g(x)}{1+b^{-x}} = g(x)\cdot \frac{2+b^x+b^{-x}}{2+b^x+b^{-x}} =g(x)\\
    \int_{-a}^a \frac{g(x)}{1+b^x}\, dx= \int_0^a \frac{g(x)}{1+b^x}+\frac{g(x)}{1+b^{-x}} \, dx = \int_0^a  g(x)\, dx
\end{gather*}

\subproblem{b}
\begin{gather*}
    I_n =\int_0^\pi \frac{\sin nx}{\sin x}\, dx\\
    I_{n+2} - I_n = \int_0^\pi \frac{\sin((n+2)x) - \sin(nx)}{\sin x} \, dx\\
    =\int_0^\pi 2\cos \paren{(n+1)x}\, dx\\
    = \frac{2}{n+1} \sin \paren{(n+1)x} \Big|^\pi_0 = 0\\
    \boxed{I_{n+2} =I_n}\\
    I_1= \int_0^\pi 1\, dx =\pi \\
    I_2 = \int_0^\pi 2\cos x\, dx = 0
\end{gather*}

\subproblem{c}
\begin{gather*}
    \int_{-\pi }^\pi \frac{\sin (2027x)}{(1+2^x)\sin x}\, dx = \int_0^\pi \frac{\sin 2027 x}{\sin x} = \boxed{\pi} \\
    \int_{-\pi }^\pi \frac{\sin (2026x)}{(1+2^x)\sin x}\, dx = \int_0^\pi \frac{\sin 2026 x}{\sin x} = \boxed{0}
\end{gather*}

\subproblem{d}
The bases cancel out upon combining two intervals together for an even function.













\end{document}

